\section{Method}
An integrated experimental-computational methodology was developed to optimise
 the properties of the poly(glycerol citrate itaconate) thermoset, supported
 by custom software enabling human-in-the-loop optimisation in a laboratory
 setting.
Programs were written in the Python programming language and developed into
 modules, whereafter these were united into a single package, ``BO4BT'', 
 available from a GitHub repository (https://github.com/ThomasDodd97/BO4BT)
 alongside its dependencies.
A .yaml configuration file is also provided which allows an identical
 Conda environment to be set up, allowing a user to run this code
 \autocite{HojlundDodd2024}.

% ========================================================================================================================================================================================
\subsection{Chemical Background}
The condensation reaction occurring between carboxyl and hydroxyl groups 
 is an esterification reaction:
\begin{equation}
    \begin{split}
        RCOOH+H&OR'\rightarrow\\
        &RCOOR'+HOH
    \end{split}
    \label{eq:Esterification}
\end{equation}
As opposed to the simple esterification shown in
 Equation~\ref{eq:Esterification}, the desired polyester thermoset,
 poly(glycerol citrate itaconate), is the product of polyesterification
 between multi-functional monomers.

Polyesterification is a form of polycondensation reaction which takes place 
 as follows: an initial ester is formed from two difunctional monomers,
 where one must at least be a diol and the other a diacid (alternatively,
 each may possess a carboxyl and a hydroxyl group), producing a molecule
 of water and an ester with two spare functional groups- these groups can
 subsequently react with more difunctional monomers causing polymerisation
 into a linear polyester (See Figure~\ref{fig:SimpleLinearEsterification}).
When monomers with a functionality greater than two are introduced into a 
 polyesterification reaction, then side chains can form along the primary
 backbone of the growing polyester, which causes the development of a
 branching polymeric structure (See Figure~\ref{fig:SimpleReactionOverview}).
When a side chain reaches another large growing polyester molecule and
 reacts with it a crosslink is said to have developed, and when widespread
 crosslinking occurs generating a 3D network of polyester molecules a
 thermoset polyester has been created.
Figure~\ref{fig:SimpleReactionOverview} shows how the monomers employed in
 the synthesis of poly(glycerol citrate itaconate) give rise to this
 branching thermoset structure.

\begin{figure*}[htb]
    \captionsetup{format=custom,width=0.8\textwidth,skip=0pt}
    \includegraphics[width=\textwidth]{SimpleLinearEsterification}
    \captionof{figure}{
            \footnotesize{
                Archetypical polyesterification between a diol (propylene
                glycol) and a diacid (succinic acid) forming a linear
                polyester.
                Functional groups participating in the reaction are
                highlighted: hydroxyl groups blue, and carboxyl groups
                orange.
                In the products, ester groups are blue and carbonyl groups
                and water are also orange in accordance with provenance.
            }
        }
    \label{fig:SimpleLinearEsterification}
\end{figure*}

Importantly, during the polyesterification process water is generated as
 the main by-product, and hence it is assumed that virtually all mass loss
 between prepolymer and final cured thermoset would be associated with
 water evaporation.
Given that each molecule of water produced corresponds to an ester bond
 formed, mass loss can be said to be a good proxy for extent of
 polymerisation within a sample \autocite{Trenkel2008}.
Thus, this research identifies mass loss as a metric of interest and sought
 to maximise this during optimisation of poly(glycerol citrate itaconate)
 synthesis: increasing the extent of polymerisation and thereby improving
 the degree of crosslinking and associated strength of these bio-based
 thermosets.

\subsection{Experimental Method}
Synthesis of the poly(glycerol citrate itaconate) thermoset was achieved 
 via an uncatalysed polyesterification reaction (See 
 Figure~\ref{fig:SimpleReactionOverview}).
Acid catalysis was avoided to reduce corrosion damage to equipment, this
 ultimately reduced the rate at which this reaction could be carried out
 but did not prevent polycondensation \autocite{Balachandran2016}.
The monomers glycerol (99+~\unit{\percent}), citric acid (99.6~\unit{\percent}), and itaconic acid
 (99+~\unit{\percent}) used were supplied by Fisher Scientific International
 and manufactured by Acros Organics.

\begin{figure*}[htb]
    \captionsetup{format=custom,width=0.8\textwidth,skip=0pt}
    \includegraphics[width=\textwidth]{SimpleReactionOverview}
    \captionof{figure}{
        \footnotesize{
                Polyesterification between glycerol, citric acid, and
                itaconic acid forming poly(glycerol citrate itaconate).
                Functional groups participating in the reaction are
                highlighted: hydroxyl groups blue, and carboxyl groups
                orange.
                In the products, ester groups are blue and carbonyl groups
                and water are orange in accordance with provenance.
            }
        }
    \label{fig:SimpleReactionOverview}
\end{figure*}

Thermoset synthesis was carried out in two steps: firstly, a homogenous 
 prepolymer resin was prepared by stirring solid monomers (acids) together 
 with liquid monomer
 (glycerol) atop a hotplate and raising the temperature to a maximum of 
 140~\unit{\degreeCelsius} until the solids were completely dissolved.
The temperature of 140~\unit{\degreeCelsius} was chosen to accelerate dissolution
 without encouraging decarboxylation of the citric acid to itaconic
 acid \autocite{Alberts2017}.
Secondly, prepolymer was cured into the desired thermoset by transferral to
 a silicone mould before being exposed to elevated temperatures
 (90~\unit{\degreeCelsius}) for an extended period (up to a week) within
 a drying oven.

Preliminary experimentation showed that temperatures exceeding
 100~\unit{\degreeCelsius}, caused boiling of byproduct water within samples.
Gaseous water from boiling was entrapped within the curing polyester creating
 bubbles within the final structure, resulting in unpredictable densities
 of the final material.
Given the unpredictable nature of samples generated in this 
 fashion, drying oven temperatures were held below 90~\unit{\degreeCelsius}
 to avoid this uncontrollable effect upon the experiments.

As defined by Equation \ref{eq:MassLoss}, determination of mass loss
 associated with the synthesis of a thermoset sample (\(\Delta m\)) was
 obtained by measuring the difference between both the initial mass of
 prepolymer resin (\(m_{p}\)) and the final mass of cured thermoset
 (\(m_{t}\)).
\begin{equation}\label{eq:MassLoss}
    \Delta m=m_{p}-m_{t}
\end{equation}
Equation~\ref{eq:PctMassLoss} defines sample mass loss (\(\Delta m\))
 normalised into a percentage format (\(\Delta m_{\%}\)), this allowed
 inter-comparison since individual samples had slightly varying initial
 prepolymer resin masses (\(m_{p}\)).
\begin{equation}\label{eq:PctMassLoss}
    \Delta m_{\%}=\frac{\Delta m\times 100}{m_{p}}
\end{equation}
Taking into account an initial prepolymer mass 
 of 3~\unit{\gram}, a mass loss range spanning 0.2-0.6~\unit{\gram}, 
 and a mass balance repeatability of 0.01~\unit{\gram} we could 
 estimate the uncertainty with respect to \(\Delta m_{\%}\) using 
 a Monte Carlo sampling procedure and some assumed distributions - 
 this equated to 0.44~\% points with \(1\sigma=0.01\)~\% points 
 (see supplementary information in Section~\ref{subsec:NoMEP}).
Residence time within the drying oven for a sample was controlled by
 periodic measurement of sample mass: mass was plotted against time, an
 exponential decay curve fitted, and the sample removed from the oven once
 a pre-determined rate of mass loss was reached (this equated to around
 -0.028~\unit{\percent\per\hour} of the initial mass).

To provide independent evidence of network formation occurring as a
 function of the sample mass loss proxy variable, a practical 
 experiment was undertaken wherein five thermoset sample formulations 
 were tested for gel fraction, with the method adapted from Susa et al.
 (2024) \autocite{Susa2024}.
The tetrahydrofuran-insoluble fraction was ascertained by firstly 
 taking an initial mass measurement \(m_{\text{initial}}\) 
 (in grams) of a cured thermoset sample.
Secondly, this sample was exposed to room-temperature tetrahydrofuran
 for a period of 12~hours under agitation, before being dried at 
 40~\unit{\degreeCelsius} until a constant mass reading was obtained 
 \(m_{\text{dried}}\) (in grams).
Thirdly, Equation~\ref{eq:GelContent} was used to ascertain a 
 tetrahydrofuran-insoluble fraction value \(g\) (in \%) for the 
 thermoset sample considered.

\begin{equation}\label{eq:GelContent}
    g=\frac{m_{\text{dried}}}{m_{\text{initial}}}
\end{equation}

For the purposes of this particular experiment all measurements were 
 undertaken using a microbalance with a readability of 
 0.01~\unit{\micro\gram}, Monte Carlo investigation of error propagation 
 found the errors associated with both \(g\) and \(\Delta m_{\%}\) to 
 therewith to be negligible (around 0.0004~\% for both metrics).

\begin{center}
    \captionsetup{format=custom,width=0.48\textwidth,skip=0pt}
    \includegraphics[width=0.48\textwidth]{MassLossInvestigation.png}
    \captionof{figure}{
        \footnotesize{
            Overview of the relationship between the mass loss proxy
             variable and THF-insoluble fraction associated with a given
             formulation.
            Numbering indicates the 3 samples which were recoverable
             and characterised during the course of experimentation.
        }
    }
    \label{fig:MassLossInvestigation}
\end{center}

Having prepared and cured five different formulations from across the 
 parameter space investigated, these samples were exposed to the 
 tetrahydrofuran as described.
Three of these remained sufficiently intact for quantitative recovery
 of the insoluble material, across these, sample mass loss closely
 followed gel fraction increase, as presented in 
 Figure~\ref{fig:MassLossInvestigation}.
Two samples disintegrated during tetrahydrofuran exposure and could not
 be quantitatively recovered, and were therewith treated as failed
 measurements rather than zero insoluble fraction.
Despite the sample number limitation, this series of measurements supports
 the use of sample mass loss under curing as an empirical indicator of
 network formation.

% ========================================================================================================================================================================================
\subsection{Optimisation Framing}
Having selected percentage mass loss within a thermoset sample
 (\(\Delta m_{\%}\)) as a suitable target variable for optimisation,
 predictor variables were defined and constrained creating a suitable
 parameter space describing a full range of prepolymer monomeric
 compositions.

The polyesterification of monomers to polymer seen in 
 Figure~\ref{fig:SimpleReactionOverview} is presented in 
 Equation~\ref{eq:StoichiometricEquation} as a parameterised stoichiometric 
 equation linking monomeric reagents glycerol (G), itaconic acid (It), and 
 citric acid (Ci) with the product polymer poly(glycerol citrate itaconate) 
 (P(GCiIt)) and byproduct water.
Equation~\ref{eq:StoichiometricEquation} introduces three stoichiometric
 variables representing the moles of glycerol (\(n_{\text{G}}\)), citric acid
 (\(n_{\text{Ci}}\)) and itaconic acid (\(n_{\text{It}}\)), and a variable representing
 the unknown number of moles of water produced (\(v\)).
\begin{equation}\label{eq:StoichiometricEquation}
    \begin{split}
        n_{\text{G}}~\text{G}+n_{\text{Ci}}~\text{Ci}&+n_{\text{It}}~\text{It}\rightarrow\\
        &1~\text{P(GCiIt)}+v~\text{H}_{2}\text{O}
    \end{split}
\end{equation}
The polyesterification reaction can be considered to consist of two possible 
 polymerisation reactions between the triol glycerol and either the triacid
 citric acid or the diacid itaconic acid.
By setting the glycerol stoichiometric variable to unity (\(n_{\text{G}}=1\)), 
 Figure~\ref{fig:PolymerisableRange} shows how the two separate reactions
 can be considered as if they were stoichiometric ranges, where there exist 
 stoichiometric bounds between which polymerisation is possible.

\begin{figure*}[htb]
    \captionsetup{format=custom,width=0.8\textwidth,skip=0pt}
    \includegraphics[width=\textwidth]{PolymerisableRange}
    \captionof{figure}{
        \footnotesize{
            Simplified chemical representations displaying the polymerisable
            range of the two different triol-acid combinations encountered
            during poly(glycerol citrate itaconate) synthesis.
            These ranges represent the stoichiometric intervals where
            difunctional and trifunctional monomers allow network formation; outside
            these bounds gelation may not occur.
            In the legend: ``G" glycerol, ``It" itaconic acid, `Ci" citric acid.
        }
    }
    \label{fig:PolymerisableRange}
\end{figure*}

These stoichiometric bounds are used to directly constrain the stoichiometric 
 variables; Equation~\ref{eq:CitricAcidBounds} describes the lower 
 (\(n_{\text{Ci}}^{l}\)) and upper (\(n_{\text{Ci}}^{u}\)) bounds placed on the citric acid 
 stoichiometric variable, whilst Equation~\ref{eq:ItaconicAcidBounds} 
 shows those for lower (\(n_{\text{It}}^{l}\)) and upper (\(n_{\text{It}}^{u}\))
 bounds on proportion of itaconic acid.
\begin{equation}\label{eq:CitricAcidBounds}
    n_{\text{Ci}}^{l}\leq n_{\text{Ci}}\leq n_{\text{Ci}}^{u}=0.5\leq n_{\text{Ci}}\leq 2.0
\end{equation}
\begin{equation}\label{eq:ItaconicAcidBounds}
    n_{\text{It}}^{l}\leq n_{\text{It}}\leq n_{\text{It}}^{u}=1.0\leq n_{\text{It}}\leq 2.0
\end{equation}
In addition to the basic polymerisable bounds imposed on the two
 stoichiometric variables, a further constraint must be imposed for a
 successful optimisation: it must be possible for permutations of
 prepolymer monomeric mixtures to allow one acid to entirely dominate the
 other whilst ensuring that a prepolymer will still polymerise - this
 constraint ensures that any polymerisable stoichiometric mixture of
 monomers is possible.
To achieve this, a three-dimensional bounded domain (a.k.a. parameter
 space, \(A\)) was devised consisting of three fundamental predictor
 variables: firstly, a citric acid strength predictor variable (\(s_1\))
 influencing the citric acid stoichiometric variable.
Secondly, an itaconic acid strength predictor variable (\(s_2\)) influencing
 the itaconic acid stoichiometric variable.
Thirdly, a balancing predictor variable (\(b_1\)) capable of trading off
 the two acid systems, influencing both the citric and itaconic acid
 stoichiometric variables.
\(A\subseteq \mathbb{R}^{d}\) describes the parameter space \(A\) as 
 a multi-dimensional real of \(d\) dimensions, where \(d=3\) given our
 three predictor variables.
All three predictor variables took constrained ranges of between 0 and 1 
 (\(s_1=[0,1],s_2=[0,1],b_1=[0,1]\)), and were linked mathematically to the 
 stoichiometric variables (via equations
 \ref{eq:CitricAcidStoichiometricPredictorLink}
 and \ref{eq:ItaconicAcidStoichiometricPredictorLink}).
\begin{equation}\label{eq:CitricAcidStoichiometricPredictorLink}
    n_{\text{Ci}}=(n_{\text{Ci}}^{l}+(n_{\text{Ci}}^{u}-n_{\text{Ci}}^{l})\times s_{1})\times b_{1}
\end{equation}
\begin{equation}\label{eq:ItaconicAcidStoichiometricPredictorLink}
    n_{\text{It}}=(n_{\text{It}}^{l}+(n_{\text{It}}^{u}-n_{\text{It}}^{l})\times s_{2})\times (1-b_{1})
\end{equation}
In summary, any random set of predictor variable values drawn from their
 respective constrained ranges of 0 to 1
 (e.g. \(s_{1}=0.75,s_{2}=0.25,b_{1}=0.5\)) could be used to calculate a
 unique permutation of stoichiometric variables (e.g. 
 \(N_{\text{Ci}}=0.81,N_{\text{It}}=0.63\)) which, on mixing of a prepolymer in
 accordance with Equation~\ref{eq:FinalStoichiometricEquation} and curing,
 would give rise to a unique polymer and associated mass loss
 (e.g. \(\Delta m_{\%}=30\)).
\begin{equation}\label{eq:FinalStoichiometricEquation}
    \begin{split}
        n_{\text{G}}~\text{G}+0\leq n_{\text{Ci}}\leq& n_{\text{Ci}}^{u}~\text{Ci}+\\
        0\leq n_{\text{It}}&\leq n_{\text{It}}^{u}~\text{It}\rightarrow \\
        1~\text{P}&\text{(GCiIt)}+v~\text{H}_{2}\text{O}
    \end{split}
\end{equation}
The sequence of calculations linking abstract predictor variables 
 (\(s_{1},s_{2},b_{1}\)) to the practical stoichiometric variables 
 (\(n_{\text{Ci}},n_{\text{It}}\)) and onwards to the target variable of mass loss 
 (\(\Delta m_{\%}\)), allows the expression of a black-box mapping between 
 predictors and targets.
This black-box mapping is described by \(f:A\rightarrow \mathbb{R}\),
 where an unknown function (\(f\)) maps any point in our parameter space 
 (\(A\)) to a single real number value (\(\mathbb{R}\)).
Equation~\ref{eq:BlackBoxFunction} more explicitly expresses the black-box 
 function in terms of our target variable and predictor variables at a 
 given sample number (\(i\)).
Note that sample number is an integer counter used to track the number of
 function evaluations carried out during a sequential optimisation.
\begin{equation}\label{eq:BlackBoxFunction}
    \Delta m_{\%}=f(\mathbf{x})=f([s_{1}^{i},s_{2}^{i},b_{1}^{i}])
\end{equation}
Our optimisation problem is thus expressed in 
 Equation~\ref{eq:OptimisingBlackBoxFunction} as the maximisation of the 
 black-box function, where the global optimum set of predictor variables 
 (\(\mathbf{x}^{*}\)) is sought.
\begin{equation}\label{eq:OptimisingBlackBoxFunction}
    \mathbf{x}^{*}=\argmaxA_{\mathbf{x}\in A} f(\mathbf{x})
\end{equation}
% ========================================================================================================================================================================================
\subsection{Optimisation Algorithm}
The sequential optimisation algorithm employed during this work is summarised
 by the flow diagram in Figure~\ref{fig:OptimisationAlgorithm}.
Successful conclusion of the optimisation is controlled by an if-else clause 
 comparing an initially specified permissible number of function evaluations 
 (a.k.a. sample limit, \(p_{\text{tot}}\)) with the current sample number (\(i\)):
 when \(i<p_{\text{tot}}\) the algorithm continues.
A second if-else statement controls whether a dataset exists from which a
 Gaussian process prior might be built for Bayesian sampling. 
When available, the algorithm uses the Gaussian process to draw a
 pre-defined number of samples (\(p_{\text{iter}}\)) from the parameter space for
 evaluation.
If the prior-forming dataset is lacking, an initial quasirandom (Sobol)
 sampling of the parameter space with a pre-defined number of deterministic
 samples (\(p_{\text{init}}\)) is initiated.
Samples drawn from the parameter space are processed into stoichiometric
 variables (\(n_{\text{Ci}},n_{\text{It}}\)), with which an experimenter can manufacture
 sample(s) before evaluating the target variables \(\Delta m_{\%}\)
 value(s) and returning these to the algorithm.
The algorithm seen in Figure~\ref{fig:OptimisationAlgorithm} was implemented
 in the Python programming language, and is incorporated into the `BO4BT' 
 package described earlier \autocite{HojlundDodd2024}.

\begin{center}
    \captionsetup{format=custom,width=0.48\textwidth,skip=0pt}
    \includegraphics[width=0.48\textwidth]{OptimisationAlgorithm}
    \captionof{figure}{
        \footnotesize{
            Flow diagram depicting the sequential optimisation algorithm
            developed.
        }
    }
    \label{fig:OptimisationAlgorithm}
\end{center}

On calculation of stoichiometric variables during an iteration, the
 algorithm merges all new sets of predictors variable values into a growing
 predictor-matrix (\(\mathbf{X}\)), shown in Equation~\ref{eq:MatrixOfPredictors}.
\begin{equation}\label{eq:MatrixOfPredictors}
    \begin{split}
        \mathbf{X}=&[\mathbf{x}^{i=1},\mathbf{x}^{i=2},...]\\
        =&[[s_{1}^{i=1},s_{2}^{i=1},b_{1}^{i=1}],[s_{1}^{i=2},...],...]
    \end{split}
\end{equation}
On conclusion of an iteration, the algorithm merges all data associated
 with samples evaluated into an optimisation-governing dataset.
Equation~\ref{eq:OptimisationDataset} presents this dataset (\(\mathbb{D}\))
 as a matrix which consists of the predictor-matrix (\(\mathbf{X}\)) and
 an array of target variable
 values (\(\mathbf{\Delta m_{\%}}\)).
\begin{equation}\label{eq:OptimisationDataset}
    \begin{split}
        \mathbb{D}=&[\mathbf{X},f(\mathbf{X})]\\
        =&[\mathbf{X},\mathbf{\Delta m_{\%}}]\\
        =&[[\mathbf{x}^{i=1},\mathbf{x}^{i=2},...],[\Delta m_{\%}^{i=1},...]]
    \end{split}
\end{equation}
% ========================================================================================================================================================================================
\subsection{Sampling Procedures}
During a sequential optimisation two forms of sampling were employed: 
 firstly, a deterministic space-filling sampling of a quasirandom nature was
 used to representatively sample the parameter space when initialising the
 optimisation.
Subsequently, a Bayesian sampling technique was used to sample-efficiently
 guide the search for the global optimum \(x^{*}\).
\subsubsection{Initial Quasirandom Sampling}
Optimisation routines were initiated by representatively sampling across
 the parameter space \(A\); achieved using the quasirandom Sobol sampling
 algorithm to draw a preset number (\(p_{\text{init}}\)) of samples from
 across \(A\).
Sobol sampling was employed as it provides the same or better
 space-filling properties across \(\mathbb{R}^{1\leq d\leq 40}\) domains
 as compared with pseudorandom or grid techniques \autocite{Garud2017}.
On evaluation of the initial representative set of samples
 (\([\mathbf{x}^{i=1},...,\mathbf{x}^{i=p_{\text{init}}}]\)) for the target
 variable values (\([\Delta m_{\%}^{i=1},...,\Delta m_{\%}^{i=p_{\text{init}}}]\)),
 an initial prior-forming dataset (\(\mathbb{D}^{i=p_{\text{init}}}\)) is produced.
In practice, only one quasirandomly sampled initial prior-forming dataset 
 (\(\mathbb{D}^{i=8}\)) was generated and used (visualised in Figure 
 \ref{fig:BayesianPriorsAbbreviated}).
As an aside, a relatively smooth topography witnessed during preliminary
 experimentation in lower dimensions was the justification for the
 use of only 8 samples during initial quasirandom sampling.

\begin{center}
    \captionsetup{format=custom,width=0.48\textwidth,skip=0pt}
    \includegraphics[width=0.48\textwidth]{BayesianPriorsAbbreviated}
    \captionof{figure}{
        \footnotesize{
            Initial 8 samples quasirandomly from across the parameter 
            space, \(A\), comprised of \(s_{1}\),
            \(s_{2}\), and \(b_{1}\) (as in equations
            \ref{eq:CitricAcidStoichiometricPredictorLink} and
            \ref{eq:ItaconicAcidStoichiometricPredictorLink}), which provide
            the prior-initialising dataset for the Bayesian routine.
            The coloration of points denotes distance from viewer with darker 
            points appearing further away.
        }
    }
    \label{fig:BayesianPriorsAbbreviated}
\end{center}

\subsubsection{Sequential Bayesian Sampling}
% Pedagogical Information which could be of use in the thesis introductory
% material:

% Sequential sampling used in this research used Bayes theorem
%  (see Equation~\ref{eq:Bayesianism}), where a posterior probability \(P(H|E)\)
%  (updated probability of hypothesis H given new evidence \(E\)) is equal to the 
%  prior probability (original probability of hypothesis) multiplied by the 
%  likelihood \(P(E|H)\) (probability of evidence given hypothesis) divided by the 
%  marginal probability \(P(E)\) (probability of evidence).
% As new evidence becomes available, the posterior probability is recycled into 
%  the equation as an updated prior alongside the new evidence and used to produce
%  a new posterior \autocite{Downey2013}.
% \begin{equation}\label{eq:Bayesianism}
%     \overset{\mathrm{posterior}}{P(H|E)}=\frac{\overset{\mathrm{prior}}{P(H)}\times \overset{\mathrm{likelihood}}{P(E|H)}}{\underset{\mathrm{marginal}}{P(E)}}
% \end{equation}
% Bayesian optimisation uses a Bayesian statistical model called a Gaussian process 
%  regression model (a form of Bayesian posterior distribution) derived from Gaussian 
%  process regression of a dataset \(\mathbb{D}\) \autocite{Rasmussen2005}.
% Equation~\ref{eq:GaussianProcess} shows how a gaussian process, 
%  \(\mathcal{GP}\),  is a distribution over functions which is completely
%  described by both a mean function, \(m\), and a covariance function, \(K\).
 
% \begin{equation}\label{eq:GaussianProcess}
%     f(\mathbf{x})\sim \mathcal{GP}(m(\mathbf{x}),K(\mathbf{x},\mathbf{x}'))
% \end{equation}

% More intuitively, Equation~\ref{eq:SurrogateFunction} describes the gaussian
%  process as a surrogate function \(\hat{f}(\cdot)\) describing the link 
%  between a candidate array of predictor values \(\mathbf{x}\) and a tuple
%  containing the mean and the variance of a normal distribution over values
%  of \(f\) at \(x\).

% \begin{equation}\label{eq:SurrogateFunction}
%     \hat{f}(\cdot):\mathbf{x}\mapsto [m(\mathbf{x}),K(\mathbf{x},\mathbf{x}')]
% \end{equation}

Sequential sampling by means of Bayesian optimisation was undertaken using
 a Gaussian process regression model, which was derived from the regression
 of an optimisation dataset \(\mathbb{D}\).
Equation~\ref{eq:GaussianProcess} shows how the Gaussian process,
 \(\mathcal{GP}\), is a distribution over functions which is completely
 described by both a mean function, \(m\), and a covariance function, \(K\).
\begin{equation}\label{eq:GaussianProcess}
    f(\mathbf{x})\sim \mathcal{GP}(m(\mathbf{x}),K(\mathbf{x},\mathbf{x}'))
\end{equation}

As an infinite-dimensional generalisation of multivariate normal
 distributions the Gaussian process inherits the properties of a conjugate
 prior \autocite{Deisenroth2020}, allowing the initial posterior
 (e.g. \(\mathcal{GP}^{i=8}\)) to be recycled as a prior alongside new
 data (e.g. \(\mathbb{D}+[\mathbf{x}^{i=9},\Delta m^{i=9}]\)) to generate
 an updated posterior (e.g. \(\mathcal{GP}^{i=9}\)) in accordance
 with Bayes' theorem.
The surrogate model evolves during sequential optimisation as new data 
 are incorporated at each iteration, \(i\).

During the fitting of the Gaussian process, a covariance function is used to
 encode beliefs regarding the families of functions likely to fit the
 data observed.
In this work we used a Mat\'ern covariance function containing a 
hyperparameter controlling smoothness of functions drawn from the 
Gaussian process.
This was set to \(\nu=2.5\), balancing the need
 for smoothness in sample paths against the drawbacks of higher 
 values of \(\nu\) (e.g. unphysical nature of infinitely 
 differentiable functions or mathematical instabilities when
 undertaking likelihood calculations).

After generating a posterior, a simple one-step lookahead sampling
 policy was used to select a suitable next point \(\mathbf{x}\) within the
 parameter space \(A\) to sample.
In this work, we focused on policies making use of the `simple reward'
 utility function (see Equation \ref{eq:UtilityFunction}) which takes the
 current dataset \(\mathbb{D}\), and from the set of sample coordinates
 evaluated so far \(X\) returns the predicted mean value \(\Delta m_{\%}\)
 of the point \(\mathbf{x}\) exhibiting the highest predicted mean
 (\(\mathrm{max}\ m(X)\)) according to the current posterior \(\mathcal{GP}\)
 \autocite{Garnett2023}.

\begin{equation}\label{eq:UtilityFunction}
    u(\mathbb{D})=\mathrm{max}\ m(X)
\end{equation}

In adopting this utility function, two different sampling policies were
 developed and investigated.
The first was predicated on the `probability of improvement' acquisition 
 function (defined in Equation~\ref{eq:ProbabilityOfImprovement}), where any 
 given putative \(\mathbf{x}\) is scored by the probability that it will 
 produce an updated dataset \(\mathbb{D}^{\prime}\) which generates a higher
 evaluation from the utility function than some absolute threshold \(\tau\).
The probability of a sample exceeding the current best is deduced using the
 current posterior \(\mathcal{GP}\)'s mean (\(m(\textbf{x})\)) and standard
 deviation (\(\sqrt{K(\mathbf{x},\mathbf{x})}\)) measured at that sample
 position.
% Note, this work made the absolute threshold equal to an evaluation of the
%  utility function using the current dataset (\(\tau=u(\mathbb{D})\)),
%  eschewing the added complication of deciding upon a marginal threshold
%  parameter \(\varepsilon\) which could have been used
%  (\(\tau=u(\mathbb{D})+\varepsilon\)) to offset any overly exploitative
%  behaviour induced by rewarding any improvement no matter the scale.
Here, the improvement threshold was defined as the current best observed
 value (\(\tau=u(\mathbb{D})\)), rather than introducing an additional
 offset parameter \(\varepsilon\).
This simplified the optimisation while avoiding unnecessary hyperparameter
 tuning.

\footnotesize
\begin{equation}\label{eq:ProbabilityOfImprovement}
    \begin{split}
        \alpha_{PI}(\mathbf{x})=&Pr(u(\mathbb{D}^{\prime})>\tau)\\
        =&\int_{\tau}^{\infty}\mathcal{N}(m(\mathbf{x}),\sqrt{K(\mathbf{x},\mathbf{x})})dm(\mathbf{x})\\
        =&1-\Phi((\tau-m(\mathbf{x}))/\sqrt{K(\mathbf{x},\mathbf{x})})
    \end{split}
\end{equation}
\normalsize

The second policy made use of the `expected improvement' acquisition
 function (defined in Equation \ref{eq:ExpectedImprovement}), which
 incorporates an expectation of improvement value
 (\(c=u(\mathbb{D}^{\prime})-u(\mathbb{D})\)).
This counteracts any overtly exploitative behaviour by rewarding both the scale and 
 probability of an improvement above the threshold \(\tau\) \autocite{Teh2020}.

\footnotesize
\begin{equation}\label{eq:ExpectedImprovement}
    \alpha_{EI}(\mathbf{x})=\int_{\tau}^{\infty}c\ \mathcal{N}(m(\mathbf{x}),\sqrt{K(\mathbf{x},\mathbf{x})})dm(x)
\end{equation}
\normalsize

Equation \ref{eq:AcquisitionFunctionOptimisation} shows how the next point
 to be sampled, \(\mathbf{x}^{+}\), is found by trivially optimising
 the cheap acquisition function (many putative samples \(\textbf{x}\) may be
 evaluated) to find the global optimum within the bounds of the parameter 
 space \(A\) \autocite{Teh2020}.

\begin{equation}\label{eq:AcquisitionFunctionOptimisation}
    \mathbf{x}^{+}=\argmaxA_{\mathbf{x}\in A}\alpha(\mathbf{x})
\end{equation}

In the present work both of these acquisition functions
 are employed during computational experiments concerning an emulated
 parameter space, whilst only probability of improvement was deployed in
 practice.
Gaussian processes and acquisition functions were implemented using two
 open-source packages: scikit-optimize ("skopt", v0.9.0) \autocite{Kernc2021}
 and Ax ("ax-platform", v0.4.3) \autocite{Facebook2024}.
Note that the Gaussian process regression models developed by these software
 packages assumed that measured percentage mass loss values comprised both
 an underlying composition-dependent response plus some independent and
 normally distributed noise.
Both the Ax and scikit-optimize packages inferred the homoscedastic noise
 during the fitting of the Gaussian process, with the \(\sigma\) value
 being equal to 0.46\% points for the scikit-optimize fitted \(\mathcal{GP}\)
 (at the culmination of the practically-undertaken optimisation routine).

\subsubsection{One-Shot Space-Filling Sampling}
To facilitate the assessment of the Bayesian techniques' sample efficacy,
 two simple one-shot sampling techniques were employed so that benchmark 
 values could be obtained.
This was achieved by using either pseudorandom or quasirandom (Sobol) sampling 
 techniques to disperse the maximum allowable number of samples 
 (\(p_{\text{tot}}=27\)) across the parameter space \(A\).
Subsequently, for any technique deployed, all samples, \(\mathbf{X}\), were 
 evaluated, \(f(\mathbf{X})\), to retrieve associated target variable 
 values, \(\mathbf{\Delta m_{\%}}\), generating a complete dataset \(\mathbb{D}\).
The highest target variable value was identified within each technique's dataset 
 (\(\Delta m_{\%}^{\circ}\)) and equated to a benchmark value.
The present work deployed both during computational experiments, whilst only
 the pseudorandom sampler saw practical application. 
Figure~\ref{fig:Benchmarks} visualises the spread of samples it generated, 
 and on functional evaluation of this set the best mass loss discovered was 
 recovered \(\Delta m_{\%}^{\circ}=14.01\).

\begin{center}
    \captionsetup{format=custom,width=0.48\textwidth,skip=0pt}
    \includegraphics[width=0.48\textwidth]{Benchmarks.png}
    \captionof{figure}{
        \footnotesize{
            Spread of the 27 samples drawn by the space-filling one-shot 
             pseudorandom sampling  technique during practical experimentation.
            The coloration of points denotes distance from viewer with darker 
             points appearing further away.
        }
    }
    \label{fig:Benchmarks}
\end{center}

\subsubsection{Sequential Space-Filling Sampling}
To provide a more robust comparator of the Bayesian techniques, a second 
 sequential sampling technique was employed during computational 
 experimentation.
Here, we re-constructed a sequential space-filling sampling technique
 developed by Crombecq et al. (2011) called the
 mc-intersite-proj-th (MIPT) algorithm \autocite{Crombecq2011}.
At each iteration during a sequential sampling routine this 
 objective-indifferent technique undertakes a Monte Carlo
 sampling exercise, wherein each of the vast number of putative candidates 
 have their ability to fill space (within the parameter space and between
 the samples evaluated so far) quantified and the strongest candidate
 selected for evaluation.

% ========================================================================================================================================================================================
\subsection{Emulation of Physical System}
Since repeated experimental optimisation was impractical, a computational
 emulator of the physical system was constructed to facilitate statistical
 comparison of the sampling strategies.
In fact, five emulators were developed, four being based on Gaussian processes
 (with differing covariance functions and hyperparameterisations) whilst the 
 fifth employed gradient boosted decision trees (GBDT).
The dataset informing the development of these emulators was comprised of all 
 experimentally evaluated samples from both preliminary investigation of
 the parameter space, preparation of the pseudorandom benchmark, and 
 the Bayesian optimisation routine employing the probability of improvement
 acquisition function.
This consolidated dataset consisted of 83 stoichiometric variable arrays
 (e.g. \({n_{ci},n_{it}}\)) and associated target variable evaluations
 (e.g. \(\Delta m_{\%}\)).
Predictor variable arrays were eschewed as the intent was to train emulators
 capable of predicting \(\Delta m_{\%}\) from \({n_{ci},n_{it}}\), reasoning
 being two-fold: firstly, predictor variable arrays are trivially converted
 into stoichiometric variable arrays.
Secondly, emulators are more easily trained in lower dimensional settings.

Subsequently, the consolidated dataset was split (in the ratio 1:2) 
 into both test and train sets.
To ensure a suitably representative train set and difficult-to-predict
 test set, a two step process was undertaken.
Firstly, Monte Carlo sampling was used to identify candidate test/train
 splits which exhibited the least spatial overlap between test and train
 samples.
Secondly, from amongst the highest scoring of candidates, a test/train
 split was manually selected which displayed both a representative train
 set spread and a test set which included at least one point in the vicinity
 of the highest \(\Delta m_{\%}\) sample.
Figure~\ref{fig:TestTrainSplit} provides an overview of the test/train
 split finally adopted.

\begin{center}
    \captionsetup{format=custom,width=0.48\textwidth,skip=0pt}
    \includegraphics[width=0.48\textwidth]{TestTrainSplit.png}
    \captionof{figure}{
        \footnotesize{
            Overview of the test/train split undertaken with
             respect to the consolidated dataset.
            Bounds describe the outer edges of the parameter space.
        }
    }
    \label{fig:TestTrainSplit}
\end{center}

Using the train set, five models were trained: firstly, a series of
 three Gaussian processes which employed the Matérn kernel were trained,
 with each of these employing a monotonically increasing value of
 \(\nu=[0.5,1.5,2.5]\) (respectively denoted as ``GP-Mat-05'', 
 ``GP-Mat-15'', and ``GP-Mat-25'').
Secondly, a separate Gaussian process was trained which employed the
 radial basis function as its kernel (denoted ``GP-Rbf'').
Thirdly, the XGBoost package was used to train a gradient-boosted
 decision tree model (denoted ``Xgb'') \autocite{ChenT2026}.
Figure~\ref{fig:EmulatedSpaces} shows a comparison of the
 five emulators with respect to the functional landscape which
 they produce when queried over the entire parameter space.

\begin{figure*}[htb]
    \captionsetup{format=custom,width=0.8\textwidth,skip=0pt}
    \includegraphics[width=\textwidth]{EmulatedSpaces.png}
    \captionof{figure}{
        \footnotesize{
            Contour plots describing the functional landscapes
             across the parameter space as emulated by the five
             different surrogate models.
            Bounds describe the outer edges of the parameter space.
            GBDT: gradient boosted decision tree.
        }
    }
    \label{fig:EmulatedSpaces}
\end{figure*}

By predicting on each array of stoichiometric variable values from the 
 withheld test dataset, a given emulator could provide a predicted target
 variable value array.
In turn, this array could be used to gauge an emulator's performance by
 measuring either the mean squared error (MSE) or the root mean squared 
 error (RMSE) between it and the withheld test dataset's array of target 
 variable values.
Table~\ref{tab:EmulatorPerformance} presents the performance of the five
 emulators with respect to these two metrics, as well as providing the
 emulator's global optimum target variable values.
It is clear that although the GP-Rbf emulator displays the strongest 
 performance in terms of RMSE and MSE, it fails to replicate the
 experimental global optimum value (\(15.14\% < 20.77\%\)); as opposed
 to the true global optimum associated with the latent objective function,
 this term denotes the globally optimal result given the experimental
 data available.
Meanwhile, emulator GP-Mat-05 provides both an emulated global optimum 
 approaching the experimental global optimum value 
 (\(18.57\% \approx 20.77\%\)) whilst retaining similarly low values for 
 both RMSE and MSE, making it a more robust model of the experimental data
 available.

\footnotesize
\begin{center}
    \begin{tblr}{colspec={Q[si={table-format=3.2},l]Q[si={table-format=3.2},l]Q[si={table-format=3.2},l]Q[si={table-format=3.2},l]},row{1}={guard},column{1}={guard}}
        {Emulator}&{RMSE\\(\%)}&{MSE\\(\%\(^{2}\))}&{Max\\(\%)}\\
        \hline
        GP-Mat-05   & 1.416 & 2.005 & 18.57 \\
        GP-Mat-15   & 1.435 & 2.058 & 18.27 \\
        GP-Mat-25   & 1.444 & 2.085 & 18.29 \\
        GP-Rbf      & 1.414 & 1.999 & 15.14 \\
        Xgb         & 1.486 & 2.209 & 19.14 \\
    \end{tblr}
    \captionsetup{format=custom,width=0.48\textwidth,skip=0pt}
    \captionof{table}{
        \footnotesize{
            Comparison of emulator performance in predicting
             upon the withheld test dataset.
            MSE: Mean squared error.
            RMSE: Root mean squared error.
            Max: The emulated global optimum target variable value.
        }
    }
    \label{tab:EmulatorPerformance}
\end{center}
\normalsize